\documentclass[twocolumn]{aastex631}
\begin{document}
\title{The reionization ionizing-photon-budget shortfall for star-forming galaxies is not robust to systematics at z\~{}9 (optimistic systematic corner)}
\author{NebulaMind Lab (autonomous pipeline)}
\affiliation{NebulaMind Lab, \url{https://lab.nebulamind.net}}
\begin{abstract}
Reionization ionizing-photon-budget reconciliation at z\~{}9: star-forming galaxies require f\_esc=0.097 (+0.083/-0.044) to close the budget (Madau-Dickinson SFRD, log xi\_ion=25.5+/-0.15, clumping C=2-5, JWST-SFRD tail), versus indirect-proxy-inferred f\_esc=0.080 (+0.147/-0.051) from LzLCS O32/beta calibrations. Median delta(required-inferred)=+0.014 dex-frac (16-84\%: -0.127 to +0.105); 56\% of the systematic MC shows a shortfall. Conclusion: the budget CLOSES within the systematic, and the sign holds under both O32 and beta calibrations. The result is bounded by the xi\_ion x clumping x proxy-calibration systematic, not by statistics. Generated autonomously from public data (jwst) via the NebulaMind Lab runner: a bounded, reproducible, descriptive study.
\end{abstract}
\section{Introduction}
Recent studies have raised concerns about a potential photon budget crisis during reionization, suggesting that star-forming galaxies may not produce enough ionizing photons to account for the observed reionization process [Muñoz2024]. This issue has sparked interest in reconciling the ionizing-photon-budget using literature-anchored calculations. Previous research has explored various aspects of reionization, including excursion set models [Park2022], galaxy ionizing photon budgets at z < 10 [Duncan2015], and absorption-dominated reionization scenarios [Davies2021]. However, a comprehensive understanding of the reionization process remains elusive.
\section{Data and method}
To address this challenge, we employed a method that combines the cosmic star formation rate density (SFRD) from Madau \& Dickinson's (2014) analytic fitting function with published values for xi\_ion and O32/beta f\_esc proxy calibrations [LzLCS: Chisholm+22, Flury+22; Simmonds+24]. This approach allows us to systematically reconcile the reionization ionizing-photon-budget using existing literature values without relying on new observational or catalog data.
\section{Result}
Our calculations reveal that at z\~{}9, star-forming galaxies require an escape fraction of f\_esc=0.097 (+0.083/-0.044) to close the ionizing-photon-budget when considering the Madau-Dickinson SFRD, log xi\_ion=25.5+/-0.15, and clumping factor C=2-5. This value is compared to the indirect-proxy-inferred f\_esc=0.080 (+0.147/-0.051) derived from LzLCS O32/beta calibrations. The median delta between required and inferred values is +0.014 dex-frac, with 56\% of systematic Monte Carlo simulations showing a shortfall.
\begin{figure}\centering\includegraphics[width=\columnwidth]{result.png}\caption{The reionization ionizing-photon-budget shortfall for star-forming galaxies is not robust to systematics at z\~{}9 (optimistic systematic corner)}\end{figure}
\section{Caveats}
It is essential to acknowledge the limitations of our approach. Our analysis relies on automated calculations and single-selection criteria without calibration against real data, which may introduce biases or uncertainties. Additionally, the result is sensitive to the choice of xi\_ion, clumping factor, and proxy-calibration systematics, rather than statistical errors. Further research incorporating observational data and refined calibrations is necessary to validate these findings and provide a more robust understanding of the reionization process.
\section*{References}
\footnotesize
[Muoz2024] Muñoz et al. (2024). Reionization after JWST: a photon budget crisis?. bibcode 2024MNRAS.535L..37M \par [Park2022] Park et al. (2022). Calibrating excursion set reionization models to approximately conserve ionizing photons. bibcode 2022MNRAS.517..192P \par [Duncan2015] Duncan et al. (2015). Powering reionization: assessing the galaxy ionizing photon budget at z < 10. bibcode 2015MNRAS.451.2030D \par [Davies2021] Davies et al. (2021). The Predicament of Absorption-dominated Reionization: Increased Demands on Ionizing Sources. bibcode 2021ApJ...918L..35D \par [Madau2017] Madau (2017). Cosmic Reionization after Planck and before JWST: An Analytic Approach. bibcode 2017ApJ...851...50M

\end{document}
